\chapter{Integration}
\section{Riemann Integrability over a box}



\section{Lebesgue's Characterization of Riemann integrable functions over a box}




\section{Fubini's Thm}
\begin{theorem}{Fubini's Thm}\label{Fubini's Thm}
    Conditions:\\
    (1) Let $Q = A \times B$, 其中 $A \subseteq \bR^k$, $B \subseteq \bR^l$ 都是 boxes.\\
    (2) Let $f: Q \rightarrow \bR$ be bounded.\\
    
    Statement: \\
    If $f$ Riem intable over $Q$, 那么 
    $$
    \ul{I}: x \mapsto \underline{\int_B} f(x,y) dy
    $$
    以及 
    $$
    \ol{I}:  x \mapsto \overline{\int_B} f(x,y) dy
    $$
    都是 \textbf{Riem intble over A} 的，并且有 

    $$
    \int_Q f = \int_A \underline{\int_B} f(x,y) \,dy \,dx = \int_A \overline{\int_B} f(x,y) \,dy \,dx
    $$
    
\end{theorem}
\begin{proof}
    一个 partition $P$ 一定可以分为 $P = (P_A, P_B)$, 它的每个 subbox $R = R_A \times R_B$.\\
    对于一个 $R$, 任取 $x_0 \in R_A$, 我们都有 $$m_R(f) \leq m_{R_B} f(x_0, \cdot)$$\\
    Summing over all $R_B \in P_B$, 我们有: 任取 $x_0 \in R_A$, 
    $$
    \sum_{R_B \in P_B} m_{R_A \times R_B} (f) v(R_B) \leq \sum_{R_B \in P_B} m_{R_B}(f(x_0, \cdot)) v(R_B) = L(f(x,\cdot), P_B) \leq I(x)
    $$
    Taking the inf over all possible $x_0$,
    $$
     \sum_{R_B \in P_B} m_{R_A \times R_B} f v(R_B) \leq m_{R_A}(\ul{I})
    $$
    而后 sum over all possible $R_A$, 得到:
    $$
    L(f,P) \leq L(\ul{I}, P_A)
    $$
    于是我们可以把 $L(\ul{I}, P_A), U(\ul{I}, P_A) $ bound 在 $L(f,P), U(f,P)$ 之间. Have four similar arguments. This gives the result.
\end{proof}



\subsection{Corollaries}

\begin{corollary}{Corollary 1}
    Conditions:和 Fubini Thm 相同\\
    Statement:\\
    $\int_B f(x,y) dy \;$  exists a.e. on $A$.
\end{corollary}
\begin{remark}
    即 $m (\{ x |  y \mapsto f(x,y) \text{ not intble over B} \}) = 0$ 
\end{remark}



\begin{corollary}{Corollary 2}{Natural Corollary}
Conditions: $Q \subseteq \bR ^ n$ 是一个 box.

Statement: 如果 $f: Q \rightarrow \bR$ ctn, 那么
$$
\int_Q f = \int^{b_1}_{a_1} \int^{b_2}_{a_2} \cdots \int^{b_n}_{a_n} f \; d x_n \cdots d x_1
$$
\end{corollary}

\begin{proof}
    平凡推论. ctn 函数, 不论 $Q$ 如何分维度 boxes, 都是 intble 的. 所以 by induction 可以一个一个维度积分.
\end{proof}






\section{Riemann Integrability over bounded sets}

\begin{definition}{integral over bounded sets}
    Conditions:\\
    (1) $S \sub \bR^n$ bdd, \\
    (2) $f: S \rightarrow \bR$ bdd\\
    Define: 
    $$
    f_S (x) = \begin{cases}
        f(x) \; \text{, if } x \in S \\
        0 \; \text{, else}
    \end{cases}
   $$
   Now we define \textbf{the integral of $f$ over $S$}: 
   $$
   \int_S f(x) dx = \int_Q f_S (x) dx
   $$
   p.t that $\int_Q f_S (x) dx$ exists.
\end{definition}
\begin{remark}
    我们把 integral 推广到任意 bounded set 上, 通过取一个 $S$ 的 covering box, 并取一个辅助函数, 在 $S$ 上面是 $f$, 在 $S$ 外面是 $0$.\\
    我们容易提出这一问题: 这个定义成立的前提是, 我们必须确定, 取不同的 covering boxes 得到的这个新函数的积分都是相同的. 下面这个 Lemma 表明了这一点. 
\end{remark}

\subsection{well-definedness of Riem integral over bdd set}
\begin{lemma}{The reason why our def of integral over bounded sets is fine:}
    Conditions:\\
    (1) $Q, Q' \sub \bR^n$ bdd.\\
    (2) $f: \bR^n \rar \bR$ supported only on $Q \intsec Q'$.\\
    Statement:\\
    (1) $f$ intble over $Q$ $\eq$ $f$ intble over $Q'$ \\
    并且 (2)
    $$
    \int_Q f = \int_{Q'} f 
    $$
\end{lemma}
\begin{proof}
    It suffices to prove it for $Q \sub Q'$. 因而只要证明了这一 case, 我们对于 $Q \not \sub Q'$ 的 case 也可以选取 $Q''$ s.t. $Q, Q' \sub Q$, 得到 $\int_Q f = \int_{Q''}f$, $\int_{Q'} f = \int_{Q''} f$.\\
    因而易证. partition 也可以通过包含构造. 多出的分割点无所谓, 因为在上面 m(f) 都是 0.
\end{proof}


\subsection{properties of Riem integral over bdd sets}
\begin{theorem}{properties of Riem integral over bdd sets}
    Conditions:\\
    (1) $S \sub \bR^n$ be bdd.\\
    (2) $f,g: S\rar \bR$ intble.\\
    Statement:
    \begin{enumerate}
        \item \textbf{Linearity of integral and intbility}: 
        $\forall c \in \bR$, $f + cg$ is also intble, 并且 
        $$ \int_S (f + cg) = \int_S f + c \int_S g$$
        \item 
        $M,m: S \rar \BR $ defined by 
        $$
        M: x \mapsto \max{(f(x), g(x))}
        $$
        $$
        m: x \mapsto \min{(f(x), g(x))}
        $$
        也是 intble 的.

        \item 
        $$f \leq g \; \forall x \implies \int_S f \leq \int_S g$$

        \item 
        $|f|$ 也是 intble 的, 并且 
        $$
        |\int_S f | \leq \int_S |f|
        $$

        \item \textbf{Monotonicity}: Let $T sub S$. If $f$ nonneg 且在 $T$ 上也 intble. Then:
        $$
        \int_T f \leq \int_S f
        $$

        \item \textbf{Additivity}: 
        $f$ intble on $S_1, S_2$ $\implies$ $f$ intble on $S_1 \union S_2$, $S_1 \intsec S_2$.\\
        并且有:
        $$
        \int_{S_1 \union S_2} = \int_{S_1} f + \int_{S_2} f - \int_{S_1 \intsec S_2} f
        $$

        \item 对于 \textbf{bdd 且 disjoint} $S_1, \cdots, S_m$, 有:
        $$
        \int_{\union_n S_n} f = \sum_n \int_{S_n} f
        $$
    \end{enumerate}
    
\end{theorem}


\subsection{Riem intble on bdd set 的等价条件}
\begin{remark}
    我们知道, $f$ intble on $S$ 的条件其实即存在一个 $Q \supset S$ 使得 $m(D_{f_S}) = 0$.\\
    我们可以把它等价为更加直接的条件:
\end{remark}

\begin{theorem}{Characterization of Riem Intbility on bdd sets}\label{Characterization of Riem Intbility on bdd sets}
    Conditions:\\
    (1) $S \sub \bR^n$ be bdd.\\
    (2) $f,g: S\rar \bR$ intble.\\
    Define:\\
    (1) $D_f = \{x \in S | f \text{ not ctn at } x\}$\\
    (2) $E = \{ x_0 \in \partial S  | \displaystyle \lim_{x \rar x_0} f(x) \not = 0 \}$
    \\Then:
    $$
    f \text{ Riem intble } \displaystyle\eq m(D_f) = m(E) = 0
    $$
\end{theorem}
\begin{remark}
    不需要证明. 实际上, 取一个 covering box, 创造出来的新函数的不连续点集比起原来的不连续点集合, 多出来的点的就是这里的 $E$.
    $$
    D_{f_S} \backslash D_f = E
    $$
\end{remark}



\subsection{Jordan measure 和 integral 的关系}
\begin{theorem}{Jordan measurable 对于 integrability 的意义} \label{equivalent conditions for Jordan-mesurability}
    Condition: $S \sub \BR^n $ is bdd.\\
    Statement: TFAE.
    \begin{enumerate}
        \item $S$ is Jordan measurable.
        \item const function $f = 1$ 在 $S$ 上 Riem intble.
        \item  $m(\partial S) = 0$
        \item  $\ol{m_J}(\partial S) = 0 $
    \end{enumerate}
\end{theorem}
\begin{proof}
    1,3,4 等价 by IBL, 显然.\\
    证明 2, 3 等价: 考虑 indicator function of S, on 一个任意的包含 S 的 box.
\end{proof}

\begin{remark}
Define: 
$$
\chi_S (x) = \begin{cases}
    1 \; , x \in S \\
    0 \; , x \not\in S
\end{cases}
$$
注意到
$$
D_{\chi_S} = \partial S
$$
\end{remark}


\begin{corollary}{To compute Jordan measure}
    $S$ Jordan measurable $\implies$ 
    $$
    m_J (S) = \int_S 1 
    $$
\end{corollary}



\begin{corollary}
    $S$ is Jordan measurable $\eq$ all ctn function from $S \rar \bR$ are Riem intble on $S$.
\end{corollary}




\section{Extended Integrals}
\begin{definition}{All Jordan measurable sets on $\bR$}
    记录一些 notations.\\
    $$
    \cJ = \{ \text{all Jordan measurable sets on } \bR^n \}
    $$
    $$
    \cJ_c = \{ \text{all Jordan measurable compact sets on } \bR^n \}
    $$
\end{definition}


\begin{definition}{negative and positive parts of a function}
    $$
    f_+ (x) = \max(f(x), 0)
    $$
    $$
    f_- (x) = \max(-f(x), 0)
    $$
\end{definition}

\begin{note}
    $f_+, f_-$ are all non-negative.\\
    并且有
    $$
    f = f_+ - f_-
    $$
    $$
    |f| = f_+ + f_-
    $$
\end{note}


\begin{remark}
    如果 $f$ ctn, 那么一定有 $f_+, f_-$ 也 ctn.
\end{remark}


\begin{definition}{extended integral}
    Condition:\\
    (1) $A \sub \BR^d$ open.\\
    (2) $f:A\rar \bR $ ctn.\\
    Define: the extended integral of $f$ over $A$:
    Case 1, for non-neg $f$:
    $$
    \int_A f = \sup_{D \sub A, D\in \cJ_c} \int_D f
    $$
    Case 2, for normal $f$:
    $$
    \int_A f = \int_A f_+ + \int_A f_-
    $$
\end{definition}
\begin{remark}
    extended integral 的理念: 对于任意一个 open set 上的 real-valued 函数, 我们把它的积分定义为这个 open set 包含的所有 compact J-measurable 子集上的函数积分的上确界.\\\\
    这个 extended def 主要服务于定义域 unbounded 的函数. 我们的 ordinary integral 只 define 在了 bounded set 上, 而我们现在把它扩展到某些 bounded set 上.\\
    我们施加的额外条件是这个函数连续. 所以现在我们对任何 $\BR^n \rar \bR$ 的连续函数都可以尝试求其 extended integral 是否存在.\\\\
    而注意到两个问题:\\
    (1) 这个定义对于 bounded set 上的连续函数也成立. 因而对于 bounded set 上的连续函数, 我们现在有两个积分定义, $ord \int_A f$ 和 $ext \int_A f$. \textbf{我们需要 justify 这两个定义对于这样的函数是等价的.}\\
    (2) 这个 extended integral 不仅 extend 了积分定义 from bounded set to unbounded set, 并且也扩展了 integrability. 即: \textbf{有些函数可以不是 Riemann integrable 的, 但是 extended integral 却存在.}
\end{remark}

\begin{example}
    下面是一个不 Riemann integrable 但是却存在 extended integral 的函数.
    考虑 $A$ bdd, open, not J-measurable(rectifiable).\\
    于是 by \ref{equivalent conditions for Jordan-mesurability}, 我们知道 $f(x) = 1, x\in A$ 是 not Riem intble 的.\\
    但是, by def of extended integral, 它却存在 extended integral.
\end{example}

\paragraph*{}现在我们 justify 这两个定义对于这样的函数是等价的. 关键在于: 任何 bounded open set 都可以用 compact elementary set 来无限逼近.

\begin{lemma}{open set can be approximated by cpt set}
    任意 open $A \in \bR^n$, 都有 a seq of compact elementary sets $(C_i)$ s.t. 
    $$
    C_i \sub C_{i+1}^{\circ} \forall i
    $$
    并且
    $$
    A = \bigunion_{i=1}^\infty C_i
    $$
\end{lemma}

\begin{remark}
    我们在 IBL 还有过另外一个 statement: 一个 open set in $\bR^n$ 都是一个 ctbl union of almost disjoint boxes. 这两个 statement 都给出了用 ctbl collection of cpt sets/ boxes 来分解 open set 的方法.
\end{remark}

\begin{proof}
    
\end{proof}






